\documentclass{article}
\usepackage[UTF8]{ctex}
\usepackage{amsmath,amssymb,amsfonts}
\usepackage{enumitem}
\usepackage{geometry}
\geometry{a4paper, margin=1in}
% === 错题本排版优化（不参与题目标识） ===
\setlength{\parindent}{0pt}
\setlength{\parskip}{0.35em plus 0.10em minus 0.08em}
\setlength{\emergencystretch}{3em}
\setlist{itemsep=0.25em, topsep=0.40em, parsep=0pt}
\allowdisplaybreaks[2]
% === 错题本排版优化结束 ===
\begin{document}
\section*{长安大学2025年数学分析真题}
\section*{185．长安大学2025年数学分析试卷}

\subsection{题目1}
1．作废。

\subsection{题目2}
2 ．求数列极限 \(\lim _{n \rightarrow \infty} n\left[\left(1+\frac{1}{n}\right)^{n}-e\right]\) 。

\subsection{题目3}
3．设 \(p \geq 0, q>0, s=p+q\) ，求级数 \(\sum_{n=1}^{\infty} \frac{1}{(s n-p)(s n+q)}\) 的和。

\subsection{题目4}
4．令 \(f(x, y)=\left\{\begin{array}{ll}\frac{1}{1+(1+x y)^{\frac{1}{\nu}}}, & 0 \leq x \leq 1,0<y \leq 1 ; \\ \frac{1}{1+e^{x}}, & 0 \leq x \leq 1, y=0 .\end{array}\right.\) 求极限 \(\lim _{y \rightarrow 0} \int_{0}^{1} f(x, y) \mathrm{d} x\) ．

\subsection{题目5}
5．应用曲面面积公式，计算半径为 \(R\) 的球体表面积。

\subsection{题目6}
6．应用积分号下求积分，计算 \(I=\int_{0}^{1} \frac{x^{b}-x^{a}}{\ln x} \mathrm{~d} x, b>a>0\) ．

\subsection{题目7}
1．求曲线 \(\left\{\begin{array}{l}x=a(\cos t+t \sin t) ; \\ y=a(\sin t-t \cos t) .\end{array}\right.\) 上任一点的法线到原点的距离．

\subsection{题目8}
2．设 \(f_{n}(x)=\frac{x}{1+n^{2} x^{2}}\) ，讨论函数序列 \(\left\{f_{n}(x)\right\}\) 和 \(\left\{f_{n}^{\prime}(x)\right\}\) 在 \((-\infty,+\infty)\) 上的一致收敛性．

\subsection{题目9}
3．计算

\[
\iint_{\sum} x \mathrm{~d} y \mathrm{~d} z+2 y \mathrm{~d} z \mathrm{~d} x+z^{2} \mathrm{~d} x \mathrm{~d} y .
\]

其中 \(\sum\) 为椭球面 \(\frac{x^{2}}{a^{2}}+\frac{y^{2}}{b^{2}}+\frac{z^{2}}{c^{2}}=1\) 的内侧，\(a, b, c\) 均为正数．

\subsection{题目10}
4．计算

\[
I=\oint_{C} \frac{e^{y}}{x^{2}+y^{2}}[(x \sin x+y \cos x) \mathrm{d} x+(y \sin x-x \cos x) \mathrm{d} y]
\]

其中 \(C: x^{2}+y^{2}=1\) ，沿逆时针方向。

\subsection{题目11}
1．设 \(\lim _{n \rightarrow \infty} a_{n}=a\) ，证明：

\begin{itemize}
  \item[（1）] \(\lim _{n \rightarrow \infty} \frac{\left[n a_{n}\right]}{n}=a\) ．
  \item[（2）] 若 \(a>0, a_{n}>0\) ，则 \(\lim _{n \rightarrow \infty} \sqrt[n]{a_{n}}=1\) ．
\end{itemize}

\subsection{题目12}
2．证明不等式 \(\frac{\tan x}{x}>\frac{x}{\sin x}, x \in\left(0, \frac{\pi}{2}\right)\) ．

\subsection{题目13}
3．设 \(f(x)\) 为定义在 \([a,+\infty)\) 上的增函数。证明： \(\lim _{x \rightarrow+\infty} f(x)\) 存在的充要条件是 \(f(x)\) 在 \([a,+\infty)\) 上有上界。

\subsection{题目14}
4．证明：利用变量代换 \(u=x^{2}-y^{2}, v=2 x y\) 可将方程 \(\frac{\partial^{2} z}{\partial x^{2}}+\frac{\partial^{2} z}{\partial y^{2}}=0\) 化为 \(\frac{\partial^{2} u}{\partial u^{2}}+\frac{\partial^{2} z}{\partial v^{2}}=0\) ．

\subsection{题目15}
5．若二元函数 \(f(x, y)\) 在点 \((0,0)\) 处连续．

\begin{itemize}
  \item[（1）] 证明：若极限 \(\lim _{(x, y) \rightarrow(0,0)} \frac{f(x, y)}{x^{2}+y^{2}}\) 存在，则函数 \(f(x, y)\) 在点 \((0,0)\) 处可微．
  \item[（2）] 举例说明，反之不成立．
\end{itemize}

四．证明题．共15分．

\begin{itemize}
  \item[（1）] 若 \(f\) 和 \(g\) 在 \([a, b]\) 上可积，则
\[
\left(\int_{a}^{b} f(x) g(x) \mathrm{d} x\right)^{2} \leq \int_{a}^{b} f^{2}(x) \mathrm{d} x \int_{a}^{b} g^{2}(x) \mathrm{d} x
\]
  \item[（2）] 若 \(f\) 在 \([a, b]\) 上可积，则
\[
\left(\int_{a}^{b} f(x) \mathrm{d} x\right)^{2} \leq(b-a) \int_{a}^{b} f^{2}(x) \mathrm{d} x
\]
\end{itemize}

五．证明题．共15分．设函数 \(f(x)\) 在 \([0,1]\) 上连续，在 \((0,1)\) 内可微，且 \(f(0)=0, f(1)=1\) ．证明：

\begin{itemize}
  \item[（1）] 存在 \(\xi \in(0,1)\) ，使得 \(f(\xi)=1-\xi\) 。
  \item[（2）] 存在不同的两点 \(\eta, \xi \in(0,1)\) ，使得 \(f^{\prime}(\eta) f^{\prime}(\xi)=1\) ．
\end{itemize}
\end{document}
